% =============================================================================
%  InsaCloud — Fiche technique : la sécurité
%  Compilation :  tectonic -X compile securite.tex
% =============================================================================
\documentclass[a4paper,11pt,oneside]{report}
% =============================================================================
%  Préambule du guide technique InsaCloud
% =============================================================================
\usepackage{fontspec}
% Noms de fichiers plutôt que noms système : la police vient du paquet
% tectonic, le document se recompile donc à l'identique sur n'importe quel poste.
\setmainfont{lmroman10-regular.otf}[
  BoldFont=lmroman10-bold.otf, ItalicFont=lmroman10-italic.otf,
  BoldItalicFont=lmroman10-bolditalic.otf]
\setsansfont{lmsans10-regular.otf}[
  BoldFont=lmsans10-bold.otf, ItalicFont=lmsans10-oblique.otf]
\setmonofont{lmmono10-regular.otf}[
  BoldFont=lmmonolt10-bold.otf, ItalicFont=lmmono10-italic.otf, Scale=0.92]
\usepackage[french]{babel}
\usepackage{graphicx,xcolor,booktabs,array,tabularx,longtable,multirow}
\usepackage{tikz}
\usetikzlibrary{positioning,arrows.meta,shapes.geometric,calc,fit,backgrounds,
                decorations.pathreplacing,patterns,shadows.blur}
\usepackage{fancyhdr,enumitem,microtype}
\usepackage[colorlinks=true,linkcolor=insared,urlcolor=insared,citecolor=insared]{hyperref}
\usepackage{caption}
\usepackage{listings}
\usepackage{tcolorbox}
\tcbuselibrary{skins,breakable}
\usepackage{titlesec}
\usepackage{emptypage}

% --- Couleurs ----------------------------------------------------------------
\definecolor{insared}{RGB}{200,16,46}
\definecolor{grisfond}{RGB}{247,248,250}
\definecolor{grisligne}{RGB}{216,220,226}
\definecolor{gristexte}{RGB}{91,97,107}
\definecolor{vertok}{RGB}{26,127,75}
\definecolor{vertfond}{RGB}{238,247,241}
\definecolor{orangepartiel}{RGB}{178,106,5}
\definecolor{orangefond}{RGB}{254,246,231}
\definecolor{rougeko}{RGB}{198,40,40}
\definecolor{bleufond}{RGB}{238,243,250}
\definecolor{bleutrait}{RGB}{43,91,158}

% --- Mise en page ------------------------------------------------------------
\captionsetup{font=small,labelfont={bf,color=insared}}
\setlist[itemize]{leftmargin=1.3em,itemsep=2pt,topsep=3pt,parsep=0pt}
\setlist[enumerate]{leftmargin=1.5em,itemsep=2pt,topsep=3pt,parsep=0pt}
\setlist[description]{leftmargin=1.3em,itemsep=3pt,topsep=3pt,parsep=0pt}

% LaTeX refuse par defaut qu'une figure occupe plus de 70 % d'une page et
% exige 20 % de texte : avec beaucoup de schemas, cela repousse les flottants
% et laisse des pages a moitie vides. On relache ces seuils.
\renewcommand{\topfraction}{0.92}
\renewcommand{\bottomfraction}{0.85}
\renewcommand{\textfraction}{0.06}
\renewcommand{\floatpagefraction}{0.75}
\setcounter{topnumber}{3}
\setcounter{bottomnumber}{2}
\setcounter{totalnumber}{4}

\pagestyle{fancy}
\fancyhf{}
\fancyhead[L]{\small\textcolor{gristexte}{\leftmark}}
\fancyhead[R]{\small\textcolor{gristexte}{InsaCloud}}
\fancyfoot[C]{\small\thepage}
\renewcommand{\headrulewidth}{0.4pt}
\renewcommand{\headrule}{\hbox to\headwidth{\color{grisligne}\leaders\hrule height \headrulewidth\hfill}}

\titleformat{\chapter}[display]
  {\normalfont\huge\bfseries}
  {\textcolor{insared}{\normalsize\MakeUppercase{\chaptertitlename} \thechapter}}
  {6pt}{\Huge}
\titlespacing*{\chapter}{0pt}{-20pt}{24pt}

% --- Extraits de code --------------------------------------------------------
\lstdefinestyle{insacloud}{
  basicstyle=\ttfamily\footnotesize,
  backgroundcolor=\color{grisfond},
  frame=single, rulecolor=\color{grisligne},
  breaklines=true, breakatwhitespace=false,
  postbreak=\mbox{\textcolor{gristexte}{$\hookrightarrow$}\space},
  showstringspaces=false,
  keywordstyle=\color{insared}\bfseries,
  commentstyle=\color{gristexte}\itshape,
  stringstyle=\color{vertok},
  numberstyle=\tiny\color{gristexte},
  xleftmargin=6pt, framexleftmargin=6pt, aboveskip=8pt, belowskip=8pt,
  extendedchars=true, inputencoding=utf8,
  literate=%
    {é}{{\'e}}1 {è}{{\`e}}1 {ê}{{\^e}}1 {ë}{{\"e}}1
    {à}{{\`a}}1 {â}{{\^a}}1 {ä}{{\"a}}1
    {î}{{\^i}}1 {ï}{{\"i}}1
    {ô}{{\^o}}1 {ö}{{\"o}}1
    {û}{{\^u}}1 {ù}{{\`u}}1 {ü}{{\"u}}1
    {ç}{{\c c}}1 {É}{{\'E}}1 {È}{{\`E}}1 {À}{{\`A}}1 {Ç}{{\c C}}1
    {œ}{{\oe}}1 {Œ}{{\OE}}1 {æ}{{\ae}}1
    {—}{{---}}1 {–}{{--}}1 {’}{{'}}1 {‘}{{'}}1
    {«}{{\guillemotleft}}1 {»}{{\guillemotright}}1
    {·}{{\textperiodcentered}}1 {→}{{$\rightarrow$}}1 {←}{{$\leftarrow$}}1
    {✔}{{\checkmark}}1 {°}{{\textdegree}}1 {€}{{\texteuro}}1
}
\lstset{style=insacloud}

\lstdefinelanguage{yamlplus}{
  keywords={true,false,null,name,hosts,become,roles,tasks,vars,when,loop,with_items,
            state,mode,owner,group,path,src,dest,notify,handlers,services,image,
            environment,volumes,networks,ports,depends_on,build,command,restart},
  keywordstyle=\color{insared}\bfseries,
  sensitive=false,
  comment=[l]{\#},
  morestring=[b]",
  morestring=[b]',
}

\lstdefinelanguage{dockerfile}{
  keywords={FROM,RUN,CMD,LABEL,EXPOSE,ENV,ADD,COPY,ENTRYPOINT,VOLUME,USER,
            WORKDIR,ARG,ONBUILD,STOPSIGNAL,HEALTHCHECK,SHELL,AS},
  keywordstyle=\color{insared}\bfseries,
  sensitive=true,
  comment=[l]{\#},
  morestring=[b]",
}

% --- Encadrés pédagogiques ---------------------------------------------------
\newtcolorbox{notion}[1]{
  breakable, enhanced, colback=bleufond, colframe=bleutrait,
  boxrule=0pt, leftrule=3pt, arc=2pt, left=8pt, right=8pt, top=6pt, bottom=6pt,
  fonttitle=\bfseries\small, title={Notion — #1},
  coltitle=bleutrait, titlerule=0pt, colbacktitle=bleufond,
  attach title to upper=\par\smallskip
}

\newtcolorbox{piege}[1]{
  breakable, enhanced, colback=orangefond, colframe=orangepartiel,
  boxrule=0pt, leftrule=3pt, arc=2pt, left=8pt, right=8pt, top=6pt, bottom=6pt,
  fonttitle=\bfseries\small, title={Piège — #1},
  coltitle=orangepartiel, titlerule=0pt, colbacktitle=orangefond,
  attach title to upper=\par\smallskip
}

\newtcolorbox{retenir}{
  breakable, enhanced, colback=vertfond, colframe=vertok,
  boxrule=0pt, leftrule=3pt, arc=2pt, left=8pt, right=8pt, top=6pt, bottom=6pt,
  fonttitle=\bfseries\small, title={À retenir},
  coltitle=vertok, titlerule=0pt, colbacktitle=vertfond,
  attach title to upper=\par\smallskip
}

% --- Raccourcis --------------------------------------------------------------
\newcommand{\ok}{\textcolor{vertok}{\textbf{Conforme}}}
\newcommand{\partiel}{\textcolor{orangepartiel}{\textbf{Partiel}}}
\newcommand{\plus}{\textcolor{insared}{\textbf{Au-delà}}}
\newcommand{\fichier}[1]{\textcolor{insared}{\texttt{#1}}}
\newcommand{\cmd}[1]{\texttt{#1}}

% Les paragraphes techniques enchaînent de longues suites en \texttt{} que TeX
% ne peut pas couper : sans marge de manœuvre, elles débordent dans la marge.
\emergencystretch=3em
\hbadness=6000
\tolerance=2000

% =============================================================================
%  Propre au document « sécurité » : une fiche par mesure
% =============================================================================
\usepackage[export]{adjustbox}

% Un bandeau d'ouverture pour chaque mesure.
\newcommand{\outil}[2]{%
  \vspace{10pt}%
  \noindent\begin{tcolorbox}[enhanced, colback=grisfond, colframe=insared,
    boxrule=0pt, leftrule=4pt, arc=2pt, left=10pt, right=10pt,
    top=7pt, bottom=7pt, breakable]
  {\Large\bfseries #1}\hfill{\small\color{gristexte}#2}
  \end{tcolorbox}\vspace{2pt}}

% Les quatre questions auxquelles chaque fiche de sécurité répond.
\newcommand{\mesure}[2]{\outil{#1}{#2}}
\newcommand{\menace}{\par\vspace{5pt}\noindent{\color{insared}\textbf{La menace}}\par\nopagebreak}
\newcommand{\parade}{\par\vspace{7pt}\noindent{\color{insared}\textbf{La parade}}\par\nopagebreak}
\newcommand{\danslecode}{\par\vspace{7pt}\noindent{\color{insared}\textbf{Dans le code}}\par\nopagebreak}
\newcommand{\limite}{\par\vspace{7pt}\noindent{\color{insared}\textbf{Ce que ça ne couvre pas}}\par\nopagebreak}

% Conservés pour compatibilité avec les macros héritées du document « technologies ».
\newcommand{\quoi}{\menace}
\newcommand{\pourquoi}{\parade}
\newcommand{\role}{\danslecode}
\newcommand{\apport}{\limite}

% Capture recadrée : on ne montre que la zone dont parle le texte.
\newcommand{\gros}[3][0.8]{%
  \begin{figure}[htbp]\centering
  \includegraphics[width=#1\textwidth,frame]{figures/#2}
  \caption{#3}\end{figure}}

\renewcommand{\topfraction}{0.92}
\renewcommand{\bottomfraction}{0.85}
\renewcommand{\textfraction}{0.06}
\renewcommand{\floatpagefraction}{0.75}


\begin{document}

\begin{titlepage}
\centering
\vspace*{2.2cm}

{\Huge\bfseries InsaCloud\par}
\vspace{6pt}
{\color{insared}\rule{4cm}{2pt}\par}
\vspace{16pt}
{\LARGE Fiche technique : la sécurité\par}
\vspace{10pt}
{\large Ce que la plateforme protège, contre quoi,\par}
{\large par quel mécanisme --- et ce qu'elle ne protège pas\par}

\vspace{2cm}
\begin{minipage}{0.84\textwidth}
\color{gristexte}\small
Une plateforme qui loue des machines Linux avec un accès \cmd{root} et des ports
ouverts sur le réseau est, par construction, une surface d'attaque. Ce document
prend chaque mécanisme de défense un par un et répond toujours aux mêmes quatre
questions : \textbf{quelle menace} il traite, \textbf{quelle parade} il oppose,
\textbf{où il se trouve} dans le code avec ses valeurs exactes, et
\textbf{ce qu'il ne couvre pas}.

\medskip
Cette dernière question n'est pas une concession : une défense dont on ignore
la limite donne une fausse assurance, et une fausse assurance est pire qu'une
absence de défense.
\end{minipage}

\vfill
{\color{gristexte}\small
Mohamed Aziz Baoueb \quad\textbullet\quad Youssef Farjallah
\quad\textbullet\quad Elyes Lemhanni\\[4pt]
Projet 3 --- \textit{Outils de déploiement de plateformes}\\
INSA Centre-Val de Loire, STI 4A FISE, 2026/2027 --- Dr.~Nadia Fettah\\[8pt]
\today\par}
\end{titlepage}

\pagenumbering{roman}
\setcounter{tocdepth}{1}
\tableofcontents
\cleardoublepage
\pagenumbering{arabic}

\chapter{Le modèle de menace}

Avant de lister des défenses, il faut dire ce qu'on défend et contre qui. Sans
cela on empile des mécanismes à la mode sans savoir lesquels comptent.

\section*{Ce que la plateforme expose}

InsaCloud n'est pas un site vitrine. Elle expose délibérément trois choses que
l'on ne met normalement pas sur un réseau :

\begin{itemize}[leftmargin=14pt,itemsep=3pt]
\item \textbf{Des comptes \cmd{root}}, un par machine louée, avec leur mot de
      passe, remis à des utilisateurs qu'on ne connaît pas.
\item \textbf{Une plage de ports ouverts} --- 8000 à 9000 --- derrière lesquels
      répondent des serveurs SSH, des terminaux web et des bureaux graphiques.
\item \textbf{Un démon Docker} sur chaque nœud, c'est-à-dire un service dont le
      contrôle équivaut à être \cmd{root} sur la machine hôte.
\end{itemize}

\begin{notion}{Pourquoi le démon Docker est l'actif le plus sensible}
Quiconque peut parler au démon Docker peut lancer un conteneur privilégié qui
monte le disque de l'hôte. Il n'y a pas d'escalade à trouver : c'est le
fonctionnement normal de Docker. « Accès au démon Docker » et « accès
administrateur de la machine » sont donc deux formulations de la même chose, et
toute décision d'architecture doit partir de là.
\end{notion}

\section*{Les quatre adversaires}

\begin{center}
\small
\begin{tabularx}{\textwidth}{@{}p{3.6cm}X@{}}
\toprule
\textbf{Qui} & \textbf{Ce qu'il cherche, et ce qui l'arrête} \\
\midrule
\textbf{Le robot d'Internet} &
Il balaie les ports, tente des mots de passe SSH et des dizaines de milliers de
chemins HTTP connus. C'est l'adversaire le plus \emph{fréquent} et le plus
bête. Parades : authentification par clé, limitation des tentatives, en-têtes
et politique de contenu stricts. \\
\addlinespace
\textbf{L'utilisateur curieux} &
Un étudiant légitime qui essaie de voir les machines d'un autre, de prolonger
sa location sans payer le temps, ou de sortir de son conteneur. C'est
l'adversaire le plus \emph{probable} dans un contexte pédagogique. Parades :
cloisonnement par propriétaire vérifié en base, capacités minimales dans le
conteneur, quotas. \\
\addlinespace
\textbf{Le locataire malveillant} &
Il loue une machine en règle, puis s'en sert comme tremplin : scanner le réseau
interne, atteindre le contrôleur, miner, attaquer l'extérieur. Il est
\emph{déjà \cmd{root}} dans son conteneur --- c'est le service rendu. Parades :
durcissement du conteneur, réseau dédié, limites de ressources, durée de vie
courte. \\
\addlinespace
\textbf{Celui qui lit le dépôt} &
Le code est sur GitHub. Un secret commis par erreur, même retiré ensuite, reste
dans l'historique. Parades : \cmd{gitleaks} sur tout l'historique à chaque
envoi, Ansible Vault, fichiers de clés exclus du dépôt. \\
\bottomrule
\end{tabularx}
\end{center}

\section*{Ce qui est hors périmètre, et pourquoi}

Un modèle de menace honnête nomme aussi ce qu'il abandonne.

\textbf{Un attaquant qui a déjà \cmd{root} sur le contrôleur} a gagné : il lit
la base, les clés, et pilote tous les nœuds. Aucune des mesures décrites ici ne
le gêne. La défense contre ce cas est organisationnelle (qui a accès à la
machine physique), pas logicielle.

\textbf{Une évasion de conteneur par une faille du noyau.} Les conteneurs
partagent le noyau de l'hôte ; une vulnérabilité du noyau les traverse. Le
durcissement réduit la surface (capacités retirées,
\cmd{no-new-privileges}) mais ne crée pas la frontière qu'offrirait une machine
virtuelle. C'est le compromis assumé de la conteneurisation, et l'énoncé le
tranche en imposant Docker.

\textbf{Le déni de service distribué.} Les limites de ressources par conteneur
protègent les \emph{voisins} d'un locataire gourmand ; elles ne protègent pas la
plateforme d'une attaque volumétrique venue de l'extérieur, qui relève de
l'hébergeur.

\begin{retenir}
Trois actifs à protéger : le \textbf{démon Docker} (équivalent à \cmd{root} sur
l'hôte), les \textbf{secrets} (clés, mots de passe des machines), et le
\textbf{cloisonnement entre locataires}. Tout le reste en découle.
\end{retenir}

\section*{Les dix couches de défense}

Le reste du document suit l'ordre d'une requête, du navigateur jusqu'au
conteneur :

\begin{center}
\small
\begin{tabularx}{\textwidth}{@{}p{0.6cm}p{4.2cm}X@{}}
\toprule
\textbf{\#} & \textbf{Couche} & \textbf{Ce qu'elle empêche} \\
\midrule
2 & Identité et authentification & qu'on se fasse passer pour un utilisateur \\
3 & Session et navigateur & qu'on détourne une session ouverte \\
4 & Secrets & qu'un secret traîne en clair quelque part \\
5 & Accès aux machines louées & qu'on entre dans la machine d'un autre \\
6 & Durcissement des conteneurs & qu'un locataire sorte de sa machine \\
7 & Contrôleur et nœuds & qu'on prenne la main sur un démon Docker \\
8 & Données et réseau & qu'on corrompe l'état ou qu'on écoute le trafic \\
9 & Chaîne d'intégration & qu'une faille entre dans le dépôt sans qu'on le voie \\
10 & Limites assumées & qu'on se croie protégé là où on ne l'est pas \\
\bottomrule
\end{tabularx}
\end{center}

\chapter{Couche 2 — L'identité}

Quatre mécanismes pour qu'un compte ne puisse pas être pris par quelqu'un
d'autre.

% -----------------------------------------------------------------------------
\mesure{Le hachage des mots de passe}{mesure 1 · \fichier{app.py}}

\menace
Une base de données finit par fuiter : sauvegarde oubliée, injection SQL,
disque revendu. Si elle contient les mots de passe en clair, tous les comptes
tombent --- et, comme les gens réutilisent leurs mots de passe, des comptes
ailleurs aussi.

\parade
On ne stocke jamais le mot de passe, mais son \textbf{empreinte}, calculée par
une fonction à sens unique et \emph{lente}. La lenteur est la propriété
recherchée : elle ne gêne pas l'utilisateur qui se connecte une fois, mais elle
ruine l'attaquant qui doit en tester des milliards.

\danslecode
\cmd{generate\_password\_hash} de Werkzeug, qui applique \textbf{PBKDF2-SHA256}
avec un sel aléatoire par compte et un grand nombre d'itérations. Le sel rend
les tables précalculées inutilisables : deux utilisateurs ayant le même mot de
passe ont deux empreintes différentes.

\begin{lstlisting}[language=Python]
user_id = db.create_user(username, generate_password_hash(password), email or None)
\end{lstlisting}

\begin{piege}{Comparer en temps constant, et même quand le compte n'existe pas}
Une vérification naïve répond \emph{plus vite} quand l'identifiant est inconnu
--- il n'y a pas d'empreinte à calculer. En chronométrant les réponses, on
énumère donc les comptes existants sans jamais réussir une connexion.

La parade tient en une ligne : quand l'utilisateur n'existe pas, on calcule
quand même une empreinte factice, pour que les deux cas coûtent le même temps.

\begin{lstlisting}[language=Python]
ok = user is not None and check_password_hash(user["password_hash"], password)
# meme cout quand le compte n'existe pas : on ne doit pas pouvoir les enumerer
check_password_hash(generate_password_hash("x"), password) if user is None else None
\end{lstlisting}
\end{piege}

\limite
Un mot de passe faible reste faible, même bien haché. D'où la mesure suivante.

% -----------------------------------------------------------------------------
\mesure{La politique de mot de passe}{mesure 2 · \fichier{config.py}}

\menace
Les attaques par dictionnaire ne cassent pas la cryptographie : elles essaient
\cmd{123456}, \cmd{password}, \cmd{azertyuiop}. Un seul compte faible suffit à
entrer.

\parade
Une longueur minimale et le refus d'une liste de mots de passe courants, à
l'inscription.

\danslecode
\begin{lstlisting}[language=Python]
PASSWORD_MIN_LENGTH = int(os.environ.get("INSACLOUD_PASSWORD_MIN", "10"))
COMMON_PASSWORDS = {"password", "motdepasse", "123456789", "1234567890",
                    "azertyuiop", ...}
\end{lstlisting}

Dix caractères, pas huit : la règle classique des huit caractères date d'une
époque où le matériel était mille fois plus lent.

\limite
Aucune politique n'empêche \cmd{Soutenance2026!} --- conforme, et devinable par
qui connaît le contexte. La vraie défense contre le devinage ciblé est la
suivante.

% -----------------------------------------------------------------------------
\mesure{La limitation des tentatives}{mesure 3 · \fichier{app.py}}

\menace
Avec un formulaire de connexion public, un robot teste des milliers de
combinaisons par minute. Sans limite, il finit par trouver.

\parade
Un compteur de tentatives échouées, sur \textbf{deux axes} et dans une fenêtre
glissante.

\danslecode
\begin{lstlisting}[language=Python]
LOGIN_WINDOW_MINUTES    = 15   # la fenetre glissante
LOGIN_MAX_FAILURES_USER = 5    # par compte vise
LOGIN_MAX_FAILURES_IP   = 20   # par adresse source
\end{lstlisting}

Les deux axes traitent deux attaques différentes :

\begin{itemize}[leftmargin=14pt,itemsep=2pt]
\item \textbf{par compte} --- l'attaquant s'acharne sur \cmd{admin} depuis mille
      adresses : le compteur du compte le bloque ;
\item \textbf{par adresse} --- l'attaquant essaie un mot de passe très probable
      sur mille comptes différents (\emph{password spraying}), ce qui ne
      déclenche aucun compteur par compte : le compteur par adresse le bloque.
\end{itemize}

Chaque tentative est enregistrée en base avec son horodatage et son adresse
source ; les entrées de plus de vingt-quatre heures sont purgées à chaque
connexion réussie, pour que la table ne grossisse pas indéfiniment.

\begin{retenir}
Deux compteurs valent mieux qu'un : un seul compteur par compte ne voit pas le
\emph{password spraying}, et un seul compteur par adresse ne voit pas l'attaque
distribuée. C'est la combinaison qui couvre les deux.
\end{retenir}

\limite
Un attaquant disposant de beaucoup d'adresses et visant beaucoup de comptes
passe sous les deux seuils. À l'échelle de ce projet c'est acceptable ; en
production on ajouterait un CAPTCHA après quelques échecs, ou une temporisation
croissante.

% -----------------------------------------------------------------------------
\mesure{La validation des entrées}{mesure 4 · \fichier{config.py}}

\menace
Toute donnée venue de l'utilisateur finit quelque part : dans une requête SQL,
dans un nom de conteneur, dans une page HTML, dans une ligne de commande. Si
elle n'est pas contrainte, elle y fait ce qu'elle veut.

\parade
Des expressions régulières \textbf{ancrées} --- \cmd{\^{}} et \cmd{\$} --- qui
décrivent ce qui est \emph{autorisé}, jamais ce qui est interdit. Une liste
d'interdits est toujours incomplète ; une liste d'autorisés ne l'est jamais.

\danslecode
\begin{lstlisting}[language=Python]
USERNAME_RE   = re.compile(r"^[a-zA-Z0-9][a-zA-Z0-9_.-]{2,31}$")  # compatible Docker
EMAIL_RE      = re.compile(r"^[^@\s]{1,64}@[^@\s.]+(\.[^@\s.]+)+$")
SSH_PUBKEY_RE = re.compile(r"^(ssh-ed25519|ssh-rsa|ecdsa-sha2-nistp(256|384|521)"
                           r"|sk-ssh-ed25519@openssh\.com) [A-Za-z0-9+/]+=*"
                           r"( [^\r\n]{0,128})?$")
MIN_DURATION_MINUTES = 1
MAX_DURATION_MINUTES = 120
MAX_INSTANCES_PER_USER = 3
\end{lstlisting}

Le format de l'identifiant n'est pas arbitraire : il est choisi pour être
\textbf{compatible avec les noms de conteneurs Docker}, puisqu'il s'y retrouve.
Un identifiant contenant une espace ou un point-virgule serait refusé bien
avant d'atteindre la ligne de commande.

\begin{notion}{La requête est bornée avant même d'être lue}
\cmd{MAX\_CONTENT\_LENGTH = 16 * 1024} : Flask rejette toute requête de plus de
seize kilooctets. Aucun formulaire légitime de la plateforme n'en approche ---
le plus gros est une clé publique SSH, quelques centaines d'octets. Cette
borne coupe court aux tentatives d'épuisement mémoire par un corps de requête
géant, avant que le code applicatif ne voie quoi que ce soit.
\end{notion}

\limite
La validation protège le format, pas le sens. Une durée de 120 minutes est
valide ; demander trois machines de 120 minutes toutes les deux minutes reste
possible, et c'est le quota (\cmd{MAX\_INSTANCES\_PER\_USER}) qui l'arrête.

\chapter{Couche 3 — La session et le navigateur}

L'utilisateur est authentifié. Reste à empêcher qu'on lui vole sa session, ou
qu'on la lui fasse utiliser à son insu.

% -----------------------------------------------------------------------------
\mesure{Le cookie de session}{mesure 5 · \fichier{config.py}}

\menace
Après la connexion, c'est le cookie qui prouve l'identité. Le voler, c'est
devenir l'utilisateur sans jamais connaître son mot de passe.

\parade
Quatre attributs, chacun fermant une porte différente.

\danslecode
\begin{lstlisting}[language=Python]
SESSION_COOKIE_HTTPONLY   = True             # inaccessible a JavaScript
SESSION_COOKIE_SAMESITE   = "Strict"         # jamais envoye depuis un autre site
SESSION_COOKIE_SECURE     = HTTPS            # transmis seulement en TLS
SESSION_COOKIE_NAME       = "insacloud_session"
PERMANENT_SESSION_LIFETIME = timedelta(hours=2)
\end{lstlisting}

\begin{center}
\small
\begin{tabularx}{\textwidth}{@{}p{3.4cm}X@{}}
\toprule
\textbf{Attribut} & \textbf{La porte qu'il ferme} \\
\midrule
\cmd{HttpOnly} & Un script injecté dans la page ne peut pas lire le cookie :
une faille XSS ne se transforme pas en vol de session. \\
\cmd{SameSite=Strict} & Le navigateur n'envoie pas le cookie quand la requête
vient d'un autre site. C'est la défense structurelle contre le CSRF, celle qui
agit avant même que la requête parte. \\
\cmd{Secure} & Le cookie ne circule jamais en clair. Activé dès que le site est
en HTTPS --- désactivé en développement local, sinon le cookie ne serait jamais
émis et rien ne fonctionnerait. \\
\cmd{2 heures} & Une session oubliée sur un poste de salle TP finit par
expirer. \\
\bottomrule
\end{tabularx}
\end{center}

\begin{piege}{Une session Flask est signée, pas chiffrée}
C'est la méprise la plus répandue. Le contenu du cookie est encodé en base64 :
n'importe qui peut le \emph{lire}. Ce que la clé secrète garantit, c'est qu'on
ne peut pas le \emph{modifier} sans invalider la signature.

Conséquence directe sur le code : la session ne contient que
l'\textbf{identifiant} de l'utilisateur. Jamais son rôle, jamais un jeton,
jamais un mot de passe. Tout le reste est relu en base à chaque requête.
\end{piege}

\limite
Un cookie volé reste valable jusqu'à son expiration : il n'est lié ni à
l'adresse IP ni au navigateur. Lier une session à son adresse casserait les
connexions mobiles et n'arrêterait qu'un attaquant peu motivé.

% -----------------------------------------------------------------------------
\mesure{Les jetons CSRF}{mesure 6 · tous les formulaires}

\menace
Un site tiers contient une image ou un formulaire qui pointe vers
\cmd{/instances/42/supprimer}. L'utilisateur, connecté à InsaCloud dans un
autre onglet, le visite : son navigateur joint son cookie, et la requête
part avec son identité.

\parade
Chaque formulaire porte un jeton aléatoire lié à la session ; le serveur refuse
toute requête qui modifie l'état sans ce jeton. Un site tiers ne peut pas le
deviner --- il ne peut pas lire la page qui le contient, la politique de même
origine l'en empêche.

\danslecode
\cmd{WTF\_CSRF\_ENABLED = True}, y compris en configuration de test : la
protection CSRF reste donc exercée par la suite de tests, et une régression qui
la désactiverait ferait échouer la chaîne d'intégration.

\begin{retenir}
\cmd{SameSite=Strict} et le jeton CSRF font le même travail à deux endroits
différents --- le navigateur et le serveur. On garde les deux : le premier
dépend du navigateur de l'utilisateur, le second de personne.
\end{retenir}

\limite
Le CSRF ne protège pas d'un script qui s'exécute \emph{dans} la page : celui-là
lit le jeton. C'est la politique de contenu qui l'en empêche.

% -----------------------------------------------------------------------------
\mesure{Les en-têtes de sécurité}{mesure 7 · \fichier{app.py}}

\menace
Le navigateur est permissif par défaut : il exécute les scripts qu'on lui
donne, affiche la page dans un cadre invisible, devine les types de contenu,
et transmet l'URL consultée au site suivant. Chacun de ces comportements est
une attaque possible.

\parade
Huit en-têtes posés sur \emph{chaque} réponse, par un \cmd{after\_request}
unique --- impossible d'en oublier un sur une route.

\danslecode
\begin{lstlisting}[language=Python]
@app.after_request
def security_headers(response):
    csp = ("default-src 'self'; script-src 'self'; style-src 'self'; "
           "img-src 'self' data:; font-src 'self'; connect-src 'self'; "
           "object-src 'none'; base-uri 'self'; form-action 'self'; "
           "frame-ancestors 'none'")
    if HTTPS:
        csp += "; upgrade-insecure-requests"
        response.headers["Strict-Transport-Security"] = \
            "max-age=31536000; includeSubDomains"
    response.headers["Content-Security-Policy"]   = csp
    response.headers["X-Content-Type-Options"]    = "nosniff"
    response.headers["X-Frame-Options"]           = "DENY"
    response.headers["Referrer-Policy"]           = "no-referrer"
    response.headers["Permissions-Policy"]        = \
        "camera=(), microphone=(), geolocation=(), payment=()"
    response.headers["Cross-Origin-Opener-Policy"]   = "same-origin"
    response.headers["Cross-Origin-Resource-Policy"] = "same-origin"
\end{lstlisting}

\begin{center}
\small
\begin{tabularx}{\textwidth}{@{}p{4.3cm}X@{}}
\toprule
\textbf{En-tête} & \textbf{Ce qu'il interdit} \\
\midrule
\cmd{Content-Security-Policy} & Charger un script, un style, une police ou une
image d'ailleurs que du site lui-même. \cmd{object-src 'none'} bloque les
greffons, \cmd{form-action 'self'} empêche un formulaire de poster ailleurs,
\cmd{base-uri 'self'} empêche de réécrire la base des liens relatifs. \\
\cmd{frame-ancestors 'none'} \newline \cmd{X-Frame-Options: DENY} & Afficher le
site dans un cadre sur un autre site --- la parade au \emph{clickjacking}. Les
deux en-têtes disent la même chose ; le second pour les navigateurs anciens. \\
\cmd{X-Content-Type-Options} & Deviner le type d'un fichier plutôt que de
croire l'en-tête : un fichier téléversé ne peut pas être réinterprété en
script. \\
\cmd{Referrer-Policy: no-referrer} & Transmettre au site suivant l'URL que
l'utilisateur consultait --- qui contient parfois un identifiant. \\
\cmd{Permissions-Policy} & Demander la caméra, le micro, la position ou le
paiement. La plateforme n'en a aucun usage ; autant les refuser explicitement. \\
\cmd{Strict-Transport-Security} & Revenir en HTTP pendant un an, y compris sur
les sous-domaines. Posé seulement si le site est réellement en TLS --- le poser
en HTTP local rendrait le site inaccessible. \\
\cmd{Cross-Origin-*} & Partager le contexte de navigation ou les ressources
avec une autre origine. \\
\bottomrule
\end{tabularx}
\end{center}

\begin{piege}{\cmd{default-src 'self'} interdit les CDN --- et c'est voulu}
Une politique stricte a un coût concret : on ne peut plus écrire
\cmd{<script src="https://cdn.tailwindcss.com">}. Tailwind est donc
\textbf{précompilé} et servi depuis le site.

Le gain dépasse la sécurité : la plateforme ne dépend plus d'un service
extérieur qui pourrait tomber, changer de contenu ou observer les visiteurs.
Une politique de contenu qu'on affaiblit pour faire marcher un CDN ne protège
plus de grand-chose --- autant la tenir.
\end{piege}

\limite
La CSP limite les dégâts d'une injection, elle ne l'empêche pas. La vraie
défense contre le XSS reste l'échappement automatique des gabarits Jinja2, que
la CSP double en dernier recours.

\chapter{Couche 4 — Les secrets}

La plateforme manipule quatre secrets : la clé de session, la clé du coffre, le
mot de passe de la base, et le jeton des agents. Plus un secret par machine
louée. Aucun ne doit se retrouver en clair, nulle part.

% -----------------------------------------------------------------------------
\mesure{Deux clés, jamais la même}{mesure 8 · \fichier{crypto.py}}

\menace
Une seule clé pour tout est commode : on la perd une fois, on perd tout. Si la
clé qui signe les cookies est aussi celle qui chiffre les mots de passe
\cmd{root}, alors la fuite d'une clé de session --- par nature plus exposée,
puisqu'elle sert à chaque requête --- livre l'accès à toutes les machines
louées.

\parade
Deux clés distinctes, de deux natures différentes, pour deux usages qui n'ont
rien à voir.

\danslecode
\begin{lstlisting}[language=Python]
def load_secret_key() -> str:
    """Cle secrete des sessions Flask."""
    return load_or_create_key("INSACLOUD_SECRET_KEY", ".secret_key",
                              lambda: secrets.token_hex(32))

# Cle de chiffrement des mots de passe root (distincte de la cle de session)
VAULT = Fernet(load_or_create_key("INSACLOUD_VAULT_KEY", ".vault_key",
                                  lambda: Fernet.generate_key().decode()))
\end{lstlisting}

Le commentaire en tête du module le dit explicitement : \emph{« Séparer les
deux garantit qu'un vol de la clé de session ne donne pas accès aux machines
louées. »}

\begin{notion}{D'où viennent les clés}
\cmd{load\_or\_create\_key} suit trois pistes, dans l'ordre :
\begin{enumerate}[leftmargin=16pt,itemsep=2pt]
\item la \textbf{variable d'environnement} --- c'est ainsi qu'Ansible les
      injecte depuis son Vault, en production ;
\item un \textbf{fichier local} déjà présent ;
\item à défaut, une clé \textbf{tirée au hasard}, écrite dans un fichier
      aussitôt mis en \cmd{0600}.
\end{enumerate}
Ce troisième cas est important : il n'existe \emph{aucune} valeur par défaut
inscrite dans le code. Une clé par défaut serait la même sur tous les
déploiements, donc publique.
\end{notion}

\limite
Les fichiers \cmd{.secret\_key} et \cmd{.vault\_key} vivent sur le disque du
contrôleur. Qui lit ce disque les obtient. Un gestionnaire externe
(HashiCorp Vault) les garderait hors de la machine, au prix d'un service
supplémentaire à déployer et à sécuriser lui-même.

% -----------------------------------------------------------------------------
\mesure{Le coffre : Fernet}{mesure 9 · \fichier{crypto.py}, \fichier{services.py}}

\menace
Le mot de passe \cmd{root} de chaque machine louée doit être conservé : pour
l'afficher à l'utilisateur quand il le redemande, et pour le faire tourner. En
clair dans la base, il tombe avec elle.

\parade
Chaque mot de passe est chiffré avant d'être écrit, par \textbf{Fernet} --- la
recette de la bibliothèque \cmd{cryptography} : AES-128 en mode CBC pour le
chiffrement, HMAC-SHA256 pour l'authentification, et un horodatage.

\danslecode
\begin{lstlisting}[language=Python]
def encrypt_secret(value: str) -> str:
    """Chiffre un secret avant de l'ecrire en base."""
    return VAULT.encrypt(value.encode()).decode()
\end{lstlisting}

Appelé aux trois endroits où un mot de passe naît ou change : à la création de
la machine, à la régénération demandée par l'utilisateur, et à la recréation
après la panne d'un nœud.

\begin{notion}{Pourquoi Fernet plutôt qu'AES « à la main »}
Fernet est un format \textbf{authentifié} : il ne se contente pas de chiffrer,
il signe. Un attaquant qui modifierait un octet du texte chiffré en base
provoquerait un \cmd{InvalidToken} au déchiffrement, et non un mot de passe
différent mais accepté. Chiffrer sans authentifier est une erreur classique,
qui ouvre des attaques par altération du texte chiffré ; Fernet la rend
impossible par construction.
\end{notion}

\limite
Le coffre protège la base, pas la mémoire : pendant qu'il sert une page, le
processus détient le mot de passe en clair. Et la clé du coffre est sur la même
machine que la base --- la séparation est logique, pas physique.

% -----------------------------------------------------------------------------
\mesure{La fenêtre du coffre}{mesure 10 · \fichier{app.py}}

\menace
Afficher en permanence les mots de passe \cmd{root} sur le tableau de bord,
c'est les offrir à quiconque passe derrière l'utilisateur --- ou à une session
restée ouverte.

\parade
Les mots de passe sont \textbf{masqués par défaut}. Pour les voir, l'utilisateur
doit ressaisir le mot de passe de \emph{son compte} ; le coffre s'ouvre alors
pour cinq minutes, puis se referme seul.

\danslecode
\begin{lstlisting}[language=Python]
VAULT_WINDOW_MINUTES = int(os.environ.get("INSACLOUD_VAULT_WINDOW", "5"))
\end{lstlisting}

Le déchiffrement n'a même pas lieu si la fenêtre est fermée : le gabarit ne
reçoit pas le secret, il reçoit \cmd{None}. Un secret qui n'est pas envoyé au
navigateur ne peut pas y être intercepté.

\begin{lstlisting}[language=Python]
secret = decrypt_secret(row["root_password"]) if (vault_open and running) else None
\end{lstlisting}

\limite
Cinq minutes suffisent à une capture d'écran. La fenêtre réduit l'exposition,
elle ne la supprime pas --- et c'est précisément pourquoi l'authentification
par clé publique (couche 5) est le meilleur chemin : elle rend le mot de passe
inutile.

% -----------------------------------------------------------------------------
\mesure{Ansible Vault}{mesure 11 · \fichier{3\_ansible/group\_vars/}}

\menace
Les secrets doivent arriver sur le serveur. Trois mauvaises façons de s'y
prendre : les écrire en clair dans le dépôt, les taper à la main à chaque
déploiement, ou les passer en variables d'environnement depuis un terminal ---
où ils restent dans l'historique du shell.

\parade
Un fichier chiffré en \textbf{AES-256}, versionné tel quel dans le dépôt,
qu'Ansible déchiffre au moment de jouer le playbook avec une phrase de passe
fournie à l'exécution.

\danslecode
\fichier{3\_ansible/group\_vars/insacloud/vault.yml} contient la clé de session,
la clé du coffre et le mot de passe de la base. Le rôle \cmd{webapp} vérifie
même leur \emph{forme} avant de déployer --- une clé Fernet doit faire
exactement 44 caractères --- plutôt que de laisser l'application échouer au
démarrage avec un message obscur.

La phrase de passe vit dans un fichier \cmd{.vault\_pass} en \cmd{0600}, exclu
du dépôt par le \fichier{.gitignore}.

\begin{retenir}
Le fichier chiffré est \emph{dans} le dépôt, et c'est l'intérêt : les secrets
voyagent avec le code, le déploiement reste entièrement automatisable, et il
n'y a aucune étape manuelle de copie de clés --- étape qui, en pratique, finit
toujours par être faite de travers.
\end{retenir}

\limite
La phrase de passe est le dernier maillon, et elle n'est protégée par rien
d'autre que les droits du fichier. Qui l'obtient obtient tout le reste.

% -----------------------------------------------------------------------------
\mesure{Les secrets sur le serveur}{mesure 12 · rôle \cmd{webapp}}

\menace
Un secret déposé sur le serveur peut fuir par trois chemins qu'on oublie :
l'unité systemd (lisible par tous), les journaux, et la liste des processus.

\parade
Les secrets sont déposés dans un \textbf{fichier séparé en \cmd{0600}}, lu par
systemd via \cmd{EnvironmentFile} --- jamais écrits dans l'unité elle-même, qui
est en lecture publique, et jamais passés en arguments de ligne de commande,
qui apparaissent dans \cmd{ps}.

\danslecode
Du côté des machines louées, le même soin : l'\fichier{entrypoint.sh} termine
par

\begin{lstlisting}[language=bash]
# Le mot de passe ne doit pas rester lisible dans l'environnement des processus
unset ROOT_PASSWORD SSH_PUBKEY
exec "$(command -v supervisord)" -c "$CONF"
\end{lstlisting}

Sans ce \cmd{unset}, le mot de passe \cmd{root} resterait lisible dans
\cmd{/proc/1/environ} pendant toute la vie de la machine.

\limite
\cmd{docker inspect} montre les variables d'environnement passées au
\cmd{docker run}, donc le mot de passe initial, à qui a accès au démon Docker.
Mais qui a accès au démon Docker a déjà tout (chapitre 1) : ce n'est pas une
fuite supplémentaire.

% -----------------------------------------------------------------------------
\mesure{La recherche de secrets commis}{mesure 13 · \cmd{gitleaks}}

\menace
Un secret commis par erreur puis retiré au commit suivant \textbf{reste dans
l'historique}, consultable par quiconque clone le dépôt. Le retirer de la
dernière version ne sert à rien.

\parade
\cmd{gitleaks} parcourt \textbf{tout l'historique} à chaque envoi, dans un
travail dédié de la chaîne d'intégration.

\danslecode
Le travail \cmd{secrets} du fichier \fichier{.github/workflows/ci.yml}. Le
\fichier{.gitignore} tient par ailleurs hors du dépôt les clés
(\cmd{.secret\_key}, \cmd{.vault\_key}), le fichier \cmd{.env} de la pile
compose, et la phrase de passe du Vault.

\limite
\cmd{gitleaks} reconnaît les formats connus --- clés AWS, jetons GitHub,
clés privées. Un secret de forme quelconque, comme un mot de passe de base de
données, peut lui échapper. D'où la règle, humaine : aucun secret n'entre dans
un fichier versionné, chiffré ou non.

\chapter{Couche 5 — L'accès aux machines louées}

C'est la surface la plus exposée de la plateforme : des ports ouverts, des
comptes \cmd{root}, et des utilisateurs qu'on ne connaît pas.

% -----------------------------------------------------------------------------
\mesure{Le mot de passe à usage unique}{mesure 14 · \fichier{services.py}}

\menace
Un mot de passe \cmd{root} par défaut, ou dérivé du nom de la machine, serait
le même partout : il suffirait d'en louer une pour entrer dans toutes les
autres.

\parade
Seize caractères tirés au sort par machine, affichés \textbf{une seule fois},
puis conservés chiffrés. L'utilisateur peut le régénérer à chaud sans détruire
sa machine.

\danslecode
Le mot de passe est injecté au lancement (\cmd{-e ROOT\_PASSWORD=...}), jamais
inscrit dans l'image. Le \fichier{Dockerfile} le dit explicitement :

\begin{lstlisting}[language=bash]
# ROOT_PASSWORD n'est VOLONTAIREMENT pas defini ici : une valeur inscrite dans
# l'image serait lisible par `docker inspect` et identique sur toutes les
# machines. La plateforme l'injecte au lancement.
\end{lstlisting}

Et si la plateforme oublie de le fournir, l'\fichier{entrypoint.sh} n'applique
pas de valeur de repli connue --- il en tire une au hasard :

\begin{lstlisting}[language=bash]
if [ -z "${ROOT_PASSWORD:-}" ]; then
    ROOT_PASSWORD="$(tr -dc 'A-Za-z2-9' < /dev/urandom | head -c 16)"
fi
\end{lstlisting}

\begin{retenir}
Une valeur par défaut est une valeur publique. Entre « échouer » et « retomber
sur une valeur connue », un système sûr échoue --- ou, comme ici, tire au
hasard et le journalise.
\end{retenir}

\limite
Le mot de passe transite par une page web et le serveur doit le connaître pour
l'afficher. C'est inhérent, et c'est la raison d'être de la mesure suivante.

% -----------------------------------------------------------------------------
\mesure{L'authentification par clé publique}{mesure 15 · \fichier{entrypoint.sh}}

\menace
Un port SSH ouvert sur Internet reçoit des milliers de tentatives par jour. Un
mot de passe de seize caractères y résiste --- mais la tentative reste
\emph{possible}, et le secret doit voyager pour être utilisé.

\parade
L'utilisateur dépose sa clé publique sur son compte. Les machines qu'il loue
\textbf{ensuite} la reçoivent, et leur \cmd{sshd} est reconfiguré pour refuser
purement et simplement la méthode « mot de passe ».

\danslecode
\begin{lstlisting}[language=bash]
if [ -n "$SSH_PUBKEY" ]; then
    mkdir -p /root/.ssh && chmod 700 /root/.ssh
    printf '%s\n' "$SSH_PUBKEY" > /root/.ssh/authorized_keys
    chmod 600 /root/.ssh/authorized_keys
    sed -i 's/^#\?PasswordAuthentication .*/PasswordAuthentication no/' \
        /etc/ssh/sshd_config
    sed -i 's/^#\?PermitRootLogin .*/PermitRootLogin prohibit-password/' \
        /etc/ssh/sshd_config
fi
\end{lstlisting}

Vérifié sur une machine louée, en forçant la méthode mot de passe :

\begin{lstlisting}[language=bash]
$ ssh -o PubkeyAuthentication=no -o PreferredAuthentications=password \
      root@127.0.0.1 -p 8096
root@127.0.0.1: Permission denied (publickey).
\end{lstlisting}

\begin{notion}{Le vrai gain : il n'y a plus de secret partagé}
Avec un mot de passe, le secret existe \textbf{à deux endroits} : l'utilisateur
le connaît, et le serveur doit le connaître pour le vérifier. Quiconque lit la
base --- ou vole la clé du coffre --- obtient l'accès à toutes les machines.

Avec une clé, le serveur ne détient que la \textbf{moitié publique}, qui ne
sert à rien à un attaquant. Compromettre la base de données d'InsaCloud ne
donne alors \emph{aucun} accès aux machines. Le compromis d'un système ne se
propage plus à l'autre.
\end{notion}

\limite
Trois limites, toutes à connaître :
\begin{itemize}[leftmargin=14pt,itemsep=2pt]
\item elle ne vaut que pour les machines louées \textbf{ensuite} --- la clé est
      injectée à la création du conteneur ;
\item le mot de passe \cmd{root} existe toujours sur la machine, car le
      terminal web et le bureau en ont besoin : la clé ferme la porte SSH, pas
      toutes les portes ;
\item une clé privée sans phrase de passe vaut ce que vaut la session de
      l'utilisateur sur son poste.
\end{itemize}

% -----------------------------------------------------------------------------
\mesure{Le terminal web passe par \cmd{login}}{mesure 16 · \fichier{entrypoint.sh}}

\menace
Un terminal dans le navigateur, c'est un shell exposé en HTTP. Mal protégé,
c'est un shell \cmd{root} offert à qui connaît l'URL.

\parade
\cmd{ttyd} ne lance pas un shell : il lance \cmd{/bin/login}, le programme de
connexion du système. L'utilisateur voit l'invite \cmd{login:} puis
\cmd{Password:}, exactement comme sur une console physique.

\danslecode
\begin{lstlisting}[language=bash]
# L'authentification est faite par `login` DANS le terminal (invite "login:"
# puis "Password:", comme une console), et non par une boite de dialogue HTTP.
\end{lstlisting}

Trois bénéfices, et le dernier n'est pas évident :

\begin{enumerate}[leftmargin=16pt,itemsep=2pt]
\item l'authentification passe par \textbf{PAM} : temporisation entre échecs,
      expiration de l'invite, journalisation --- gratuitement ;
\item l'expérience est familière, donc l'utilisateur n'est pas tenté de
      contourner ;
\item \textbf{couper l'authentification SSH par mot de passe n'enferme personne
      dehors} : le terminal web continue de fonctionner, puisqu'il ne passe pas
      par SSH. L'utilisateur qui perd sa clé garde une porte.
\end{enumerate}

\begin{piege}{La version précédente utilisait Basic Auth, et c'était pire}
Protégé par une authentification HTTP basique, le terminal affichait une page
blanche dans plusieurs navigateurs --- la négociation se faisait avant
l'ouverture de la WebSocket. Le remplacement par \cmd{login} a corrigé le
défaut d'usage \emph{et} amélioré la sécurité : Basic Auth transmet le secret à
chaque requête, \cmd{login} une seule fois dans un canal déjà établi.
\end{piege}

\limite
\cmd{ttyd} sert le terminal en HTTP si aucun certificat n'est monté. Sur le
parc déployé, le rôle \cmd{tls\_cert} fournit un certificat par nœud et
\cmd{ttyd} passe en \cmd{-{}-ssl} ; en pile locale de démonstration, le trafic
est en clair sur \cmd{localhost}.

% -----------------------------------------------------------------------------
\mesure{Le bureau graphique}{mesure 17 · \fichier{entrypoint.sh}}

\menace
Un serveur VNC exposé est une cible classique : le protocole est ancien et son
authentification est faible.

\parade
Trois précautions : le port VNC n'est \textbf{jamais publié}, le mot de passe
du bureau est \textbf{distinct} de celui de \cmd{root}, et le pont vers le
navigateur passe en TLS dès qu'un certificat est disponible.

\danslecode
\begin{lstlisting}[language=bash]
Xvnc :1 -geometry 1280x800 -depth 24 -rfbport 5901 -localhost \
        -SecurityTypes VncAuth -PasswordFile /root/.vnc/passwd -AlwaysShared
websockify $NOVNC_TLS --web $NOVNC_DIR 6080 localhost:5901
\end{lstlisting}

\cmd{-localhost} est le point essentiel : le port 5901 n'est joignable que
depuis l'intérieur du conteneur. Le seul port publié est le 6080, celui de
noVNC. Le mot de passe du bureau est dérivé du mot de passe \cmd{root} mais
tronqué à huit caractères, et stocké dans un fichier en \cmd{0600}.

\begin{piege}{Huit caractères, et ce n'est pas un choix}
Le protocole VncAuth \textbf{tronque le mot de passe à huit caractères} : c'est
une limite du protocole, pas un oubli. Il a donc fallu exposer un mot de passe
de bureau distinct de celui de \cmd{root}, pour ne pas laisser croire à
l'utilisateur que son mot de passe de seize caractères protège aussi son
bureau.
\end{piege}

\limite
\textbf{VncAuth est faible} : un défi-réponse DES des années 1990. C'est
acceptable ici pour trois raisons cumulées --- la machine est éphémère, le mot
de passe n'est visible qu'en ouvrant le coffre pendant cinq minutes, et le
transport est chiffré dès qu'un certificat est monté. Sur une plateforme
permanente, il faudrait un proxy authentifié devant noVNC.

\chapter{Couche 6 — Le durcissement des conteneurs}

Le locataire est \cmd{root} dans sa machine : c'est le service vendu. Toute la
question est de faire en sorte que ce \cmd{root}-là ne soit pas \cmd{root} sur
le nœud.

% -----------------------------------------------------------------------------
\mesure{Les capacités retirées}{mesure 18 · \fichier{docker\_ops.py}}

\menace
Par défaut, Docker accorde à un conteneur une quinzaine de \emph{capacités} du
noyau Linux. Certaines sont anodines ; d'autres permettent de modifier le
réseau de l'hôte, de monter des systèmes de fichiers ou de charger des modules.

\parade
On retire \textbf{tout}, puis on rend uniquement ce dont une machine Linux
ordinaire a besoin pour démarrer ses services. C'est une liste blanche, pas une
liste noire.

\danslecode
\begin{lstlisting}[language=bash]
--cap-drop ALL
--cap-add CHOWN --cap-add DAC_OVERRIDE --cap-add FOWNER --cap-add FSETID
--cap-add SETGID --cap-add SETUID --cap-add SYS_CHROOT
--cap-add NET_BIND_SERVICE --cap-add KILL --cap-add AUDIT_WRITE
\end{lstlisting}

\begin{center}
\small
\begin{tabularx}{\textwidth}{@{}p{3.3cm}X@{}}
\toprule
\textbf{Capacité gardée} & \textbf{Pourquoi elle est nécessaire} \\
\midrule
\cmd{CHOWN}, \cmd{FOWNER}, \cmd{FSETID} & Créer les répertoires de service avec
les bons propriétaires et droits (\cmd{/root/.ssh}, \cmd{/run/sshd}). \\
\cmd{DAC\_OVERRIDE} & Permettre à \cmd{root} d'écrire là où les droits ne le
prévoient pas --- comportement normal de \cmd{root}. \\
\cmd{SETGID}, \cmd{SETUID} & \cmd{login} et \cmd{sshd} changent d'identité pour
ouvrir une session : sans elles, aucune connexion. \\
\cmd{SYS\_CHROOT} & \cmd{sshd} isole son processus non privilégié par
\cmd{chroot} pendant l'authentification. \\
\cmd{NET\_BIND\_SERVICE} & Écouter sur un port inférieur à 1024 --- le 22. \\
\cmd{KILL} & \cmd{supervisord} doit pouvoir arrêter et relancer ses services. \\
\cmd{AUDIT\_WRITE} & PAM écrit un enregistrement d'audit à chaque connexion. \\
\bottomrule
\end{tabularx}
\end{center}

\begin{notion}{Ce qui est retiré, et qui compte}
\cmd{SYS\_ADMIN} (monter des systèmes de fichiers --- la capacité la plus
dangereuse, souvent qualifiée de « nouveau \cmd{root} »), \cmd{NET\_ADMIN}
(reconfigurer le réseau, poser une interface en mode espion),
\cmd{SYS\_MODULE} (charger un module noyau), \cmd{SYS\_PTRACE} (inspecter la
mémoire d'autres processus), \cmd{MKNOD} (créer un fichier de périphérique,
donc accéder au disque brut).

Chacune de ces capacités est une voie d'évasion documentée. Les retirer ne
retire rien à l'usage légitime d'une machine louée.
\end{notion}

\limite
Les capacités cloisonnent le noyau, elles ne le remplacent pas. Une
vulnérabilité du noyau lui-même traverse la frontière.

% -----------------------------------------------------------------------------
\mesure{Pas d'élévation de privilèges}{mesure 19 · \fichier{docker\_ops.py}}

\menace
Un binaire \cmd{setuid} à l'intérieur du conteneur peut rendre des privilèges
qu'on vient justement de retirer. Beaucoup d'images en contiennent sans le
savoir (\cmd{ping}, \cmd{mount}, \cmd{su}).

\parade
\cmd{-{}-security-opt no-new-privileges:true} : le noyau refuse qu'un processus
du conteneur gagne des privilèges qu'il n'avait pas au départ, quel que soit le
bit \cmd{setuid} du binaire exécuté.

\danslecode
Dans la constante \cmd{CONTAINER\_HARDENING}, appliquée à \emph{toutes} les
machines louées sans exception --- ce n'est pas une option du formulaire.

\limite
Le drapeau empêche de \emph{gagner} des privilèges. Il ne retire pas ceux que
\cmd{root} possède déjà dans le conteneur.

% -----------------------------------------------------------------------------
\mesure{Les limites de ressources}{mesure 20 · \fichier{docker\_ops.py}}

\menace
Un locataire lance une bombe à fourche (\cmd{:(){ :|:\& };:}), un mineur de
cryptomonnaie ou une compilation massive. Sans limite, il prend tout le nœud et
toutes les autres machines deviennent inutilisables. C'est le déni de service
le plus facile à monter : il suffit d'être client.

\parade
Trois plafonds, posés au lancement.

\danslecode
\begin{lstlisting}[language=bash]
--pids-limit 512            # nombre maximal de processus : arrete la bombe a fourche
--cpus 1                    # un coeur au plus
--memory 256m | 1g          # selon le mode, terminal ou bureau
--memory-reservation ...    # le minimum garanti en cas de contention
\end{lstlisting}

La distinction entre \cmd{-{}-memory} et \cmd{-{}-memory-reservation} est celle
que le cours impose : la \emph{limite} est un plafond dur que le conteneur ne
peut pas dépasser, la \emph{réservation} est le minimum qui lui reste garanti
quand le nœud est sous pression. Sans réservation, une machine bien élevée peut
être étranglée par une voisine gourmande, même si celle-ci respecte sa limite.

\begin{retenir}
\cmd{-{}-pids-limit} est la plus sous-estimée des trois : une bombe à fourche ne
consomme ni beaucoup de CPU ni beaucoup de mémoire par processus, mais épuise
la table des processus du \emph{noyau}, qui est partagée avec l'hôte. Les
limites de CPU et de mémoire ne l'arrêtent pas.
\end{retenir}

\limite
Les limites valent par conteneur. Trois machines par utilisateur, c'est trois
fois la limite --- d'où le quota.

% -----------------------------------------------------------------------------
\mesure{Un réseau dédié}{mesure 21 · \fichier{config.py}, \fichier{docker\_ops.py}}

\menace
Sur le réseau \cmd{bridge} par défaut de Docker, tous les conteneurs se voient.
Un locataire peut donc scanner et atteindre les machines des autres locataires
--- et les services de la plateforme elle-même.

\parade
Les machines louées sont placées sur un réseau Docker \textbf{propre au
déploiement}, créé à la demande.

\danslecode
\begin{lstlisting}[language=Python]
CONTAINER_NETWORK = os.environ.get("INSACLOUD_NETWORK", f"insacloud_{CLUSTER}_net")
\end{lstlisting}

\begin{piege}{\cmd{-{}-internal} aurait coupé Internet aux locataires}
La première version créait le réseau avec \cmd{-{}-internal}, qui l'isole
totalement --- y compris de l'extérieur. L'intention était bonne, l'effet
désastreux : plus aucune machine louée ne pouvait faire \cmd{apt-get install},
ce qui est précisément ce pour quoi on loue une machine.

Le défaut a été vu avant livraison, en essayant réellement d'installer un
paquet dans une machine fraîchement louée. C'est le genre d'erreur qu'aucune
relecture ne détecte, parce que le code est exactement conforme à l'intention.
\end{piege}

\limite
Le réseau dédié sépare les machines louées des services de la plateforme ; il
ne sépare pas les locataires \emph{entre eux}. Un réseau par machine serait
plus strict, au prix d'une complexité de gestion sans commune mesure avec le
gain à cette échelle.

% -----------------------------------------------------------------------------
\mesure{L'identité par étiquette}{mesure 22 · \fichier{docker\_ops.py}}

\menace
Le Faucheur détruit les conteneurs qui ne correspondent à aucune location
active. S'il les reconnaît par leur \emph{nom}, une collision de préfixe lui
fait détruire autre chose --- par exemple les services de la plateforme, qui
s'appellent \cmd{insacloud-web}, \cmd{insacloud-db}, \cmd{insacloud-faucheur}.

\parade
Trois étiquettes Docker, apposées à chaque machine, qui servent d'identité.
Jamais le nom.

\danslecode
\begin{lstlisting}[language=Python]
MACHINE_LABEL = "insacloud.role"      # = rented-machine
CLUSTER_LABEL = "insacloud.cluster"   # = ce deploiement
WORKER_LABEL  = "insacloud.worker"    # = le noeud qui l'heberge

MACHINE_FILTERS = [f"label={MACHINE_LABEL}={MACHINE_LABEL_VALUE}",
                   f"label={CLUSTER_LABEL}={CLUSTER}"]
\end{lstlisting}

Les deux premiers filtres sont \emph{combinés} : une machine doit être une
machine louée \textbf{et} appartenir à ce déploiement. Sans le second, le
Faucheur d'une plateforme détruirait les machines d'une autre tournant sur le
même démon --- ce qui est arrivé en développement.

Le troisième, \cmd{insacloud.worker}, a été ajouté plus tard : il dit sur quel
nœud la machine a été créée. Sur un parc de machines virtuelles, la frontière
des démons suffit ; dès que deux nœuds partagent un démon, elle disparaît, et
sans ce label la réconciliation de l'un détruisait les machines de l'autre.

\begin{retenir}
Ce défaut n'a pas été trouvé en relisant le code --- il a été trouvé en faisant
tourner un parc à deux nœuds. Une étiquette est une \emph{donnée portée par
l'objet lui-même} ; un nom est une convention que rien ne garantit. Pour une
opération destructive, la différence n'est pas théorique.
\end{retenir}

\limite
Un attaquant ayant accès au démon Docker peut poser lui-même ces étiquettes sur
un conteneur à lui, et le faire adopter par la plateforme. Mais, une fois de
plus : qui a le démon a déjà tout.

% -----------------------------------------------------------------------------
\mesure{La commande Docker, sans interpréteur}{mesure 23 · \fichier{docker\_ops.py}}

\menace
L'énoncé impose de lancer les machines par la commande \cmd{docker run}. Écrite
en concaténant des chaînes et passée à un shell, une telle commande est une
injection en puissance : une distribution nommée
\cmd{alpine; rm -rf /} suffirait.

\parade
La commande est construite en \textbf{liste d'arguments} et exécutée sans
shell. Aucun interpréteur ne relit jamais les valeurs.

\danslecode
\begin{lstlisting}[language=Python]
res = subprocess.run(cmd, capture_output=True, text=True, timeout=timeout)
# pas de shell=True : les arguments sont passes en liste, donc jamais
# reinterpretes par un interpreteur de commandes.
\end{lstlisting}

Et en amont, les seules valeurs d'origine utilisateur --- durée, distribution,
mode --- sont validées : la durée est un entier borné, la distribution et le
mode sont choisis dans des dictionnaires fermés. Une valeur hors de ces
ensembles n'atteint jamais \cmd{docker}.

\begin{notion}{Deux barrières plutôt qu'une}
La validation en amont suffirait. L'absence de shell suffirait. On fait les
deux, parce que la première dépend d'un développeur qui n'oublie pas un
chemin, et la seconde d'une propriété du système. Quand l'une cède, l'autre
tient.
\end{notion}

\limite
L'exécution sans shell ne protège pas d'un argument \emph{valide} mais
malveillant --- par exemple un chemin de montage arbitraire, si un jour on
laissait l'utilisateur en choisir un. La liste des arguments est entièrement
construite par la plateforme : c'est ce qui doit le rester.

\chapter{Couche 7 — Le contrôleur et les nœuds}

Le contrôleur doit créer et détruire des conteneurs sur des machines distantes.
C'est le lien le plus sensible de l'architecture : il porte, de bout en bout,
un pouvoir équivalent à \cmd{root}.

% -----------------------------------------------------------------------------
\mesure{SSH plutôt que l'API Docker}{mesure 24 · \fichier{workers.py}}

\menace
La façon évidente de piloter un démon Docker à distance est d'exposer son API
sur un port TCP. C'est aussi la plus dangereuse : cette API n'a, par défaut,
\textbf{aucune authentification}, et y accéder revient à être administrateur de
la machine. Les moteurs de recherche de services exposés en recensent des
dizaines de milliers, et ils sont compromis en quelques minutes.

\parade
Le transport SSH natif du client Docker. Aucun port supplémentaire n'est
ouvert : SSH était déjà là, et il apporte l'authentification que l'API n'a pas.

\danslecode
\begin{lstlisting}[language=Python]
def docker_command(worker: str = "local") -> list:
    """Prefixe de commande docker pour ce noeud."""
    host = worker_host(worker)
    if host is None:
        return ["docker"]
    return ["docker", "-H", f"ssh://{DOCKER_SSH_USER}@{host}"]
\end{lstlisting}

Le module le résume en tête : \emph{« Rien n'est exposé sur le réseau à part
SSH ; l'authentification se fait par clé. »}

\begin{notion}{Ce que fait réellement \cmd{docker -H ssh://}}
Le client Docker ouvre une session SSH et y exécute
\cmd{docker system dial-stdio} sur le nœud. L'API circule alors \emph{dans} le
tunnel SSH, authentifiée et chiffrée, sans qu'aucune socket réseau n'écoute.
Si l'on avait voulu sécuriser l'API TCP, il aurait fallu monter une autorité de
certification, émettre un certificat client par contrôleur et les renouveler :
beaucoup de pièces mobiles pour un résultat équivalent.
\end{notion}

\limite
Qui obtient la clé SSH du contrôleur obtient le démon Docker de tous les nœuds.
Cette clé est donc l'actif le plus sensible du parc --- d'où la mesure
suivante.

% -----------------------------------------------------------------------------
\mesure{La clé de pilotage}{mesure 25 · rôle \cmd{controller\_key}}

\menace
Une clé SSH réutilisée --- celle du développeur, par exemple --- étend sa portée
à tout le parc, et survit au départ de son propriétaire.

\parade
Une clé \textbf{dédiée au pilotage}, générée par Ansible, déposée sur les
nœuds pour un \textbf{compte de service sans mot de passe}.

\danslecode
Le rôle \cmd{controller\_key} crée le compte et la paire de clés ; le rôle
\cmd{worker} dépose la partie publique sur chaque nœud. Le compte
\cmd{insacloud} n'a pas de mot de passe --- il n'est donc joignable que par
clé, et uniquement depuis le contrôleur.

Sur le parc à deux nœuds en conteneurs, le même principe est rejoué par le
service \cmd{cles}, qui fabrique la paire une seule fois dans un volume
partagé. L'\fichier{entrypoint} du nœud installe la partie publique et, détail
qui a coûté une heure, déverrouille le compte :

\begin{lstlisting}[language=bash]
# « adduser -D » laisse le compte VERROUILLE (un « ! » dans /etc/shadow), et
# sshd refuse un compte verrouille MEME en authentification par cle. L'etoile
# dit « aucun mot de passe n'ouvre ce compte », ce qui est l'intention.
sed -i 's|^insacloud:!|insacloud:*|' /etc/shadow
\end{lstlisting}

\begin{piege}{Un compte sans mot de passe n'est pas un compte verrouillé}
Les deux états se ressemblent et n'ont rien à voir. \cmd{!} signifie
« verrouillé » : \cmd{sshd} refuse la connexion, y compris par clé.
\cmd{*} signifie « aucun mot de passe ne peut ouvrir ce compte » : la
connexion par clé fonctionne. Le symptôme était un
\cmd{Connection closed by remote host} muet côté client ; seul le journal du
serveur disait \cmd{User insacloud not allowed because account is locked}.
\end{piege}

\limite
La clé n'a pas de phrase de passe : elle doit être utilisable par un service,
sans intervention humaine. Sa protection repose donc entièrement sur les droits
du fichier et sur la sécurité du contrôleur.

% -----------------------------------------------------------------------------
\mesure{Le jeton des agents}{mesure 26 · \fichier{api.py}, \fichier{worker\_agent.py}}

\menace
L'inscription des nœuds est dynamique : un agent appelle
\cmd{POST /workers/register} et le contrôleur l'ajoute au parc. Sans contrôle,
n'importe qui déclare un nœud, reçoit des machines à héberger --- et donc les
mots de passe de leurs locataires.

\parade
Un jeton partagé, exigé dans un en-tête à l'inscription comme au battement de
cœur.

\danslecode
\begin{lstlisting}[language=Python]
"X-Agent-Token": TOKEN        # cote agent
\end{lstlisting}

Côté contrôleur, les routes \cmd{/workers/register} et \cmd{/workers/heartbeat}
refusent la requête sans jeton valide --- ce qu'un test vérifie
explicitement (\cmd{test\_api\_enregistrement\_refuse\_sans\_jeton}). Le jeton
est distribué par Ansible Vault, jamais inscrit dans le code.

L'agent, lui, refuse de démarrer sans jeton plutôt que de tenter une
inscription anonyme :

\begin{lstlisting}[language=Python]
if not TOKEN:
    log.error("INSACLOUD_AGENT_TOKEN n'est pas defini : arret.")
\end{lstlisting}

\limite
Un jeton partagé est le même pour tous les nœuds : il ne distingue pas les
agents entre eux, et le compromettre permet d'en déclarer un faux. Une
authentification mutuelle par certificat serait plus robuste ; elle demanderait
une autorité de certification que le projet n'a pas.

% -----------------------------------------------------------------------------
\mesure{Un nœud muet n'est jamais interprété}{mesure 27 · \fichier{faucheur.py}}

\menace
Le Faucheur réconcilie Docker et la base : il détruit les conteneurs inconnus
et marque arrêtées les machines sans conteneur. Si un nœud ne répond pas et
qu'il interprète ce silence comme « il n'y a plus rien là-bas », il déclare
mortes des machines qui tournent parfaitement --- et les détruira au retour du
nœud.

\parade
Un nœud injoignable est \textbf{ignoré}, pas interprété. Aucune décision n'est
prise sur son compte.

\danslecode
\begin{lstlisting}[language=Python]
if res.returncode != 0:
    log.warning("Noeud %s injoignable pour la reconciliation : %s",
                worker, res.stderr.strip()[:120])
    continue          # on ne touche a rien pour ce noeud
\end{lstlisting}

Un test dédié verrouille ce comportement :
\cmd{test\_noeud\_injoignable\_ne\_detruit\_rien} --- \emph{« Exigence de
résilience : un worker éteint ne doit pas faire perdre l'état en base. »}

\begin{retenir}
C'est une décision de sécurité autant que de fiabilité. La question posée au
système n'est pas « que s'est-il passé ? » mais « que puis-je affirmer ? ».
Un nœud silencieux ne permet d'affirmer rien du tout, et une opération
destructive exige une certitude.
\end{retenir}

\limite
La prudence a un coût : une machine réellement disparue reste marquée active
tant que son nœud ne répond pas. C'est le bon côté sur lequel se tromper.

% -----------------------------------------------------------------------------
\mesure{Le pare-feu, différencié par groupe}{mesure 28 · rôle \cmd{security\_hardening}}

\menace
Un nœud ouvre son port SSH pour le contrôleur, et des centaines de ports pour
les machines louées. Le contrôleur, lui, n'a besoin que de HTTPS. Appliquer la
même politique aux deux ouvre inutilement des portes.

\parade
Une politique par défaut en refus, puis des règles distinctes selon le groupe
Ansible auquel la machine appartient.

\danslecode
Le rôle \cmd{security\_hardening} applique, via \cmd{community.general.ufw} :
politique entrante par défaut, SSH avec \textbf{limitation anti-force-brute}
(\cmd{ufw limit}) et restriction d'origine, ports applicatifs, et la plage des
machines louées \emph{uniquement sur les workers}.

\begin{lstlisting}
- name: Autoriser SSH avec limitation anti-force-brute sur le port {{ ufw_ssh_port }}
  community.general.ufw:
    rule: limit
    port: "{{ ufw_ssh_port }}"

- name: Autoriser les plages de ports des machines louees
  community.general.ufw:
    ...
\end{lstlisting}

\begin{piege}{L'ordre des rôles n'est pas indifférent}
Le rôle passe \textbf{avant} le rôle \cmd{docker} : \cmd{ufw enable} recharge
\cmd{iptables} et effacerait les chaînes que Docker y installe. Inversé, le
pare-feu casserait le réseau des conteneurs --- panne difficile à diagnostiquer,
parce que les deux rôles réussissent.
\end{piege}

\limite
UFW filtre le trafic entrant vers l'hôte. Docker insère ses propres règles
\cmd{iptables} pour les ports publiés, qui peuvent contourner UFW selon la
configuration : la plage des machines louées est donc ouverte de fait dès que
Docker publie les ports.

\chapter{Couche 8 — Les données et le réseau}

Deux sujets qu'on range rarement sous « sécurité » et qui en relèvent
pourtant : l'intégrité de l'état, et la confidentialité du transport.

% -----------------------------------------------------------------------------
\mesure{Le cloisonnement entre locataires}{mesure 29 · \fichier{database.py}}

\menace
C'est la faille la plus banale des applications multi-utilisateurs : l'URL
\cmd{/instances/42} affiche la machine numéro 42, et personne ne vérifie
qu'elle appartient bien à celui qui la demande. On change le numéro, on voit
les machines des autres --- et leur mot de passe \cmd{root}.

\parade
Le filtrage par propriétaire est fait \textbf{dans la requête SQL}, pas dans le
gabarit ni dans une condition d'affichage. Une donnée qui n'appartient pas à
l'utilisateur ne quitte jamais la base.

\danslecode
\begin{lstlisting}[language=Python]
requete = requete.where(Rental.user_id == user_id)
...
.where(Rental.user_id == user_id, Instance.status == INSTANCE_RUNNING)
...
.where(Archive.user_id == user_id, Archive.expires_at > utc_now())
\end{lstlisting}

La condition est présente sur les instances, sur les locations \emph{et} sur
les archives --- l'archive du travail d'un autre serait tout aussi
compromettante que sa machine.

\begin{retenir}
Filtrer à l'affichage laisse la donnée arriver jusqu'au gabarit : il suffit
d'une page oubliée, d'une réponse JSON, d'un message d'erreur trop bavard pour
qu'elle sorte. Filtrer dans la requête rend la fuite structurellement
impossible. C'est le test le plus important de la suite, et il porte un nom
explicite : \emph{cloisonnement entre comptes}.
\end{retenir}

\limite
Le cloisonnement repose sur la discipline : chaque nouvelle requête doit porter
la condition. Un modèle de données où la propriété serait imposée par le moteur
(\emph{row-level security} de PostgreSQL) l'appliquerait sans dépendre du
développeur.

% -----------------------------------------------------------------------------
\mesure{Les contraintes d'intégrité}{mesure 30 · \fichier{models.py}}

\menace
Deux machines attribuées au même port sur le même nœud : la seconde ne démarre
pas, ou pire, un utilisateur se connecte à la machine d'un autre. Une
vérification faite en Python ne tient pas sous la concurrence --- deux requêtes
simultanées peuvent lire le même « port libre » avant que l'une n'écrive.

\parade
La règle est tenue par le \textbf{moteur de base de données}, par un
\textbf{index unique partiel} : le couple (nœud, port) est unique, mais
seulement parmi les machines en cours d'exécution.

\danslecode
Le « partiel » est essentiel : sans lui, un port utilisé par une machine
détruite resterait réservé pour toujours, et la plage finirait épuisée. La
contrainte ne s'applique donc qu'aux instances \cmd{running}.

Sept tests vérifient ces règles au niveau du moteur
(\fichier{tests/test\_contraintes.py}) : qu'un port en conflit lève bien une
\cmd{IntegrityError}, que le port est libéré quand l'instance expire, que
l'identifiant et l'adresse de courriel sont uniques, que les clés étrangères
interdisent une location sans utilisateur --- avec
\cmd{PRAGMA foreign\_keys = ON} pour SQLite, qui ne les applique pas par
défaut.

\begin{notion}{Pourquoi le moteur et pas le code}
Une règle tenue par le moteur reste vraie même si une requête concurrente, un
script d'administration ou un futur appel oublient de la vérifier. Une règle
tenue seulement par le code applicatif ne tient que tant que personne ne
contourne ce code. L'énoncé demande explicitement de « tester les contraintes
d'intégrité » : la distinction est exactement celle-là.
\end{notion}

\limite
Les contraintes garantissent la cohérence, pas la justesse métier. Rien
n'empêche d'attribuer un port valide mais déjà pris par un service de l'hôte
--- d'où une vérification complémentaire côté système
(\cmd{port\_is\_free\_on\_host}).

% -----------------------------------------------------------------------------
\mesure{TLS de bout en bout}{mesure 31 · rôle \cmd{reverse\_proxy}}

\menace
Le mot de passe de connexion, le cookie de session et les mots de passe
\cmd{root} circulent entre le navigateur et le serveur. En clair, ils sont
lisibles par quiconque partage le réseau --- un réseau d'école, typiquement.

\parade
nginx termine le TLS devant Gunicorn, avec des protocoles et des suites
choisis plutôt que laissés par défaut.

\danslecode
\begin{lstlisting}
listen 443 ssl http2;
ssl_protocols TLSv1.2 TLSv1.3;
ssl_prefer_server_ciphers on;
ssl_ciphers ECDHE-ECDSA-AES256-GCM-SHA384:ECDHE-ECDSA-CHACHA20-POLY1305:
            ECDHE-ECDSA-AES128-GCM-SHA256;
client_max_body_size 16k;
\end{lstlisting}

Trois décisions :

\begin{itemize}[leftmargin=14pt,itemsep=3pt]
\item \textbf{TLS 1.2 minimum} --- les versions 1.0 et 1.1 sont obsolètes et
      vulnérables ; aucun navigateur en service ne les exige.
\item \textbf{Uniquement des suites ECDHE avec AEAD} : \cmd{ECDHE} assure la
      \emph{confidentialité persistante} (un enregistrement du trafic reste
      indéchiffrable même si la clé du serveur fuite plus tard), et GCM ou
      ChaCha20-Poly1305 chiffrent \emph{et} authentifient.
\item \cmd{client\_max\_body\_size 16k} --- la même borne qu'en applicatif, mais
      appliquée par nginx : une requête démesurée est coupée avant d'atteindre
      Python.
\end{itemize}

\begin{notion}{\cmd{ProxyFix}, ou pourquoi tout le monde avait la même IP}
Derrière un proxy inverse, Flask voit l'adresse du proxy --- \cmd{127.0.0.1}
--- et non celle du visiteur. Le compteur anti-force-brute par adresse
compterait donc toutes les tentatives du monde sur une seule, et bloquerait
tout le monde au premier attaquant. \cmd{ProxyFix} (Werkzeug) lit l'en-tête
\cmd{X-Forwarded-For} posé par nginx et rétablit la vraie adresse. Une mesure
de sécurité mal câblée peut se retourner en déni de service.
\end{notion}

\limite
Le certificat est \textbf{auto-signé} : le navigateur avertit. C'est la seule
alerte que l'analyse dynamique signale encore, et c'est le premier élément des
travaux à court terme (Let's Encrypt). Un certificat auto-signé chiffre
correctement, mais ne prouve pas l'identité du serveur : il ne protège donc pas
d'un attaquant capable de s'interposer.

\chapter{Couche 9 — La chaîne DevSecOps}

Toutes les mesures précédentes se dégradent : une dépendance vieillit, un
développeur pressé affaiblit une règle, un secret se glisse dans un commit. La
chaîne d'intégration est ce qui les maintient vraies dans le temps.

\section*{Onze travaux, à chaque envoi}

\begin{center}
\small
\begin{tabularx}{\textwidth}{@{}p{3.2cm}p{2.5cm}X@{}}
\toprule
\textbf{Travail} & \textbf{Outil} & \textbf{Ce qu'il empêche d'entrer} \\
\midrule
\cmd{lint-python} & ruff & Du code mort, des imports inutilisés, des motifs
dangereux (\cmd{shell=True}, \cmd{assert} en production). Vérifie aussi que les
schémas du \cmd{README} sont à jour. \\
\cmd{tests-python} & pytest & Une régression fonctionnelle. 169 tests, dont le
cloisonnement entre comptes, le CSRF et l'anti-force-brute. \\
\cmd{sast-python} & bandit & Une faiblesse \emph{dans notre code} : secret en
dur, aléa non cryptographique, appel système non maîtrisé. \\
\cmd{sca-dependances} & pip-audit & Une dépendance portant une vulnérabilité
connue. C'est ainsi qu'a été repérée une faille de Werkzeug, corrigée en
passant de 3.1.8 à 3.1.9. \\
\cmd{secrets} & gitleaks & Un secret commis, y compris s'il a été retiré
depuis : l'outil parcourt \textbf{tout l'historique}. \\
\cmd{dast} & nuclei + \cmd{dast.sh} & Une faiblesse visible de l'extérieur sur
l'application \textbf{en train de tourner} : en-tête manquant, cookie mal
posé, route accessible sans session. \\
\cmd{lint-ansible} & ansible-lint & Un playbook non idempotent, une tâche sans
nom, un module obsolète. Profil \cmd{production}, le plus strict. \\
\cmd{images-docker} & hadolint + Trivy & Un \cmd{Dockerfile} mal écrit, et des
paquets vulnérables \emph{dans l'image livrée} --- ce que l'audit des
dépendances ne voit pas. \\
\cmd{integration} & \cmd{integration.sh} & Une plateforme qui ne démarre pas
sur une machine vierge. Les six scénarios du sujet. \\
\cmd{publier-images} & Syft + ghcr.io & Une image sans \textbf{inventaire de
composants} (SBOM) : on doit pouvoir répondre « cette image contient-elle la
bibliothèque vulnérable annoncée ce matin ? ». \\
\cmd{deployer} & Ansible & Un déploiement manuel, non reproductible. \\
\bottomrule
\end{tabularx}
\end{center}

\begin{notion}{SAST, SCA, DAST : trois angles, pas trois noms}
\textbf{SAST} lit \emph{notre} code sans l'exécuter --- il voit les erreurs que
nous écrivons. \textbf{SCA} confronte \emph{les dépendances} à une base de
vulnérabilités --- il voit les erreurs des autres, dont nous héritons.
\textbf{DAST} attaque l'application \emph{en fonctionnement} --- il voit ce
qu'aucune lecture de code ne montre : un en-tête oublié par la configuration
nginx, un cookie mal posé en production.

Les trois sont nécessaires parce qu'aucun ne voit ce que les deux autres
voient. La plupart des failles réelles tiennent à l'assemblage, pas à une
ligne.
\end{notion}

% -----------------------------------------------------------------------------
\mesure{Le seuil de blocage}{mesure 32 · \fichier{ci.yml}}

\menace
Une chaîne qui signale tout sans rien bloquer devient un bruit de fond qu'on
cesse de lire. Une chaîne qui bloque sur tout devient un obstacle qu'on
contourne.

\parade
Un seuil explicite, justifié, différent selon l'outil.

\danslecode
\cmd{hadolint} échoue au niveau \cmd{warning}, avec trois règles écartées et
chacune commentée --- par exemple l'épinglage des versions de paquets, écarté
parce que les dépôts Debian et Alpine ne conservent pas les anciennes versions :
une image épinglée cesserait de se construire à la première mise à jour du
dépôt, et c'est la version de la \emph{distribution} qui fixe déjà l'ensemble.

Trivy ignore les vulnérabilités \cmd{-{}-ignore-unfixed} : bloquer sur une faille
sans correctif disponible n'apporte rien --- on ne peut rien en faire --- et
noierait les failles réellement corrigeables.

\begin{retenir}
Un seuil bien choisi se reconnaît à ceci : chaque exception est écrite, datée
et justifiée dans un fichier versionné. Une exception non écrite est un oubli
qui se transmet.
\end{retenir}

\limite
Les seuils reflètent un jugement, qui vieillit. Une règle écartée aujourd'hui
pour de bonnes raisons peut devenir importante ; rien ne force à les relire.

% -----------------------------------------------------------------------------
\mesure{La version de Python, imposée}{mesure 33 · \fichier{ci.yml}}

\menace
Le projet passe les tests sur le poste du développeur et échoue au
déploiement, parce que la cible n'a pas la même version du langage. C'est
l'erreur la plus banale et la plus coûteuse, et elle a frappé ce projet : une
bibliothèque épinglée n'existait pas pour Python 3.10.

\parade
La chaîne force \cmd{PYTHON\_VERSION: 3.10} --- la version qu'Ubuntu 22.04
fournit, c'est-à-dire la cible réelle. Les tests tournent avec le Python de la
production, pas celui du développeur.

\limite
L'alignement porte sur la version du langage, pas sur l'ensemble du système :
seul le travail d'intégration, qui démarre réellement la pile, couvre le reste.

% -----------------------------------------------------------------------------
\mesure{Ce que la chaîne a réellement trouvé}{mesure 34 · bilan}

Une chaîne d'intégration ne vaut pas par sa configuration mais par ses prises.
Huit défauts ont été trouvés par la chaîne ou par l'exécution réelle, que la
relecture n'avait pas vus. Les trois plus instructifs :

\textbf{Un Faucheur qui détruisait sa propre plateforme.} Il reconnaissait les
machines louées par leur préfixe de nom ; les services de la pile
s'appellent \cmd{insacloud-web}, \cmd{insacloud-db}. Correction : l'identité
par étiquette (mesure 22).

\textbf{Une réconciliation qui détruisait les machines des autres nœuds.} Sur
un parc où deux nœuds partagent un démon Docker, chacun voyait les machines de
l'autre comme orphelines. Correction : l'étiquette \cmd{insacloud.worker}, et
deux tests --- dont un qui vérifie qu'une machine \emph{sans} étiquette reste
nettoyable, pour ne pas abandonner celles créées par une version antérieure.

\textbf{Une erreur 500 une fois sur trois, au démarrage seulement.} Gunicorn
tourne avec \cmd{-{}-preload} : il importe l'application \emph{avant} de forker
ses quatre workers, et l'initialisation de la base ouvre alors une connexion
PostgreSQL qui se range dans le pool du processus maître. Les quatre fils en
héritaient et se servaient de la \textbf{même socket} simultanément. Mesuré
avant correction : trois échecs sur huit premières requêtes parallèles, à
chacun des trois redémarrages d'essai ; après correction, vingt-quatre sur
vingt-quatre.

\begin{piege}{Le défaut qui ne se voit qu'au premier regard}
Ce dernier cas mérite d'être retenu pour sa forme : le message d'erreur
(\cmd{connection in transaction status INTRANS}) n'avait \emph{aucun} rapport
avec la requête en cours, et le symptôme disparaissait tout seul après
quelques appels. Un défaut qui se répare seul est un défaut qu'on croit
imaginaire --- et c'est exactement celui qu'aurait rencontré un visiteur
ouvrant la page pour la première fois.

Le test de non-régression a été validé en le faisant d'abord \textbf{échouer},
correction désactivée. Un test qu'on n'a pas vu échouer ne prouve rien.
\end{piege}

\chapter{Couche 10 — Les limites assumées}

Ce chapitre est le plus important du document. Une défense dont on ignore la
limite donne une fausse assurance, et une fausse assurance est pire qu'une
absence de défense : on arrête de surveiller.

\section*{Les risques résiduels, et pourquoi ils sont acceptés}

\begin{center}
\small
\begin{tabularx}{\textwidth}{@{}p{3.6cm}Xp{1.9cm}@{}}
\toprule
\textbf{Risque} & \textbf{Pourquoi il est accepté ici, et ce qu'il faudrait en production} &
\textbf{Gravité} \\
\midrule
\textbf{Certificat auto-signé} &
Le trafic est chiffré, mais l'identité du serveur n'est pas prouvée : un
attaquant capable de s'interposer pourrait se faire passer pour la plateforme.
Seule alerte que l'analyse dynamique signale encore. Correction prévue :
Let's Encrypt. & Moyenne \\
\addlinespace
\textbf{VncAuth} &
Défi-réponse DES, mot de passe tronqué à huit caractères par le protocole.
Acceptable parce que la machine est éphémère, que le mot de passe n'est visible
que cinq minutes et que le transport est chiffré. En production : un proxy
authentifié devant noVNC. & Moyenne \\
\addlinespace
\textbf{\cmd{docker.sock} monté dans la pile compose} &
En mode mono-hôte, le conteneur de l'application reçoit la socket Docker de
l'hôte --- donc un pouvoir équivalent à \cmd{root} sur la machine. Le fichier
\cmd{docker-compose.yml} le dit en toutes lettres : \emph{« À n'utiliser que
sur une machine de développement. »} Sur le parc déployé, ce montage
\textbf{disparaît} : le contrôleur pilote les workers par SSH. & Élevée
(développement seulement) \\
\addlinespace
\textbf{\cmd{root} dans le conteneur loué} &
C'est le service vendu, pas un défaut. Atténué par les capacités retirées,
\cmd{no-new-privileges} et les limites de ressources. Une évasion par faille du
noyau reste possible : seule une machine virtuelle créerait la vraie
frontière. & Acceptée \\
\addlinespace
\textbf{Un seul jeton pour tous les agents} &
Il ne distingue pas les nœuds entre eux ; le compromettre permet d'en déclarer
un faux, donc de recevoir des machines à héberger. Une authentification
mutuelle par certificat demanderait une autorité de certification. & Moyenne \\
\addlinespace
\textbf{La reprise ne restitue pas les fichiers} &
Elle recrée un conteneur, elle ne déplace pas un disque. Sans stockage partagé
entre nœuds, seul le service repart. C'est le premier élément des travaux à
moyen terme. & Acceptée \\
\addlinespace
\textbf{Les clés sur le disque du contrôleur} &
\cmd{.secret\_key} et \cmd{.vault\_key} sont en \cmd{0600}, mais sur la même
machine que la base. La séparation des deux clés est logique, pas physique. Un
gestionnaire externe les garderait hors de la machine. & Moyenne \\
\addlinespace
\textbf{Pas de journalisation centralisée} &
Les traces d'accès aux machines louées restent sur chaque nœud. Une
investigation après incident serait laborieuse, et un attaquant ayant pris un
nœud peut effacer ses traces. Travaux à moyen terme. & Moyenne \\
\addlinespace
\textbf{Pas de second facteur} &
Un mot de passe compromis suffit à prendre un compte. Travaux à moyen terme. &
Faible (contexte pédagogique) \\
\bottomrule
\end{tabularx}
\end{center}

\section*{Ce qui est hors de portée du projet}

Trois choses qu'aucune mesure du document ne traite, et qu'il serait malhonnête
de laisser croire couvertes :

\textbf{Un attaquant déjà administrateur du contrôleur} lit la base, les clés,
et pilote tous les nœuds. La défense est organisationnelle --- qui a accès à la
machine --- pas logicielle.

\textbf{Une vulnérabilité du noyau Linux} traverse la frontière des conteneurs,
qui partagent ce noyau avec l'hôte. Le durcissement réduit la surface, il ne
crée pas d'isolation matérielle.

\textbf{Un déni de service volumétrique} venu de l'extérieur relève de
l'hébergeur. Les limites par conteneur protègent les locataires les uns des
autres, pas la plateforme d'une attaque réseau.

\section*{Récapitulatif : menace, parade, preuve}

\begin{center}
\small
\begin{tabularx}{\textwidth}{@{}p{4.1cm}p{4.4cm}X@{}}
\toprule
\textbf{Menace} & \textbf{Parade} & \textbf{Preuve} \\
\midrule
Fuite de la base & PBKDF2 salé, Fernet &
Les mots de passe en base sont illisibles sans la clé du coffre. \\
Énumération des comptes & Empreinte factice si le compte n'existe pas &
Le temps de réponse est identique dans les deux cas. \\
Force brute sur la connexion & 5 échecs par compte, 20 par adresse, 15 min &
\cmd{test\_securite.py} \\
Vol de session & \cmd{HttpOnly}, \cmd{SameSite=Strict}, \cmd{Secure}, 2 h &
En-têtes vérifiés par l'analyse dynamique. \\
CSRF & Jeton sur tous les formulaires, actif jusqu'en test &
\cmd{test\_securite.py} \\
Injection de script & CSP \cmd{default-src 'self'}, aucun CDN &
En-tête vérifié par l'analyse dynamique. \\
Vol de la clé de session & Clé du coffre distincte &
\fichier{crypto.py} : deux clés, deux origines. \\
Secret commis dans Git & gitleaks sur tout l'historique &
Travail \cmd{secrets} de la chaîne. \\
Devinage du mot de passe d'une machine & Clé publique, \cmd{PasswordAuthentication no} &
\cmd{Permission denied (publickey)} \\
Évasion de conteneur & \cmd{cap-drop ALL}, \cmd{no-new-privileges} &
\cmd{test\_docker.py} vérifie la commande construite. \\
Épuisement du nœud & \cmd{pids-limit}, \cmd{cpus}, \cmd{memory} &
\cmd{test\_docker.py} \\
Prise de contrôle d'un démon Docker & SSH par clé dédiée, API jamais exposée &
\fichier{workers.py} : \cmd{docker -H ssh://} \\
Faux nœud déclaré & Jeton d'agent exigé &
\cmd{test\_api\_enregistrement\_refuse\_sans\_jeton} \\
Destruction sur un nœud muet & Le silence n'est pas interprété &
\cmd{test\_noeud\_injoignable\_ne\_detruit\_rien} \\
Accès aux machines d'un autre & Filtrage par propriétaire en SQL &
\cmd{test\_securite.py}, cloisonnement entre comptes \\
Collision de port & Index unique partiel en base &
\cmd{test\_contraintes.py} (7 cas) \\
Écoute du trafic & TLS 1.2/1.3, suites ECDHE + AEAD &
Configuration nginx du rôle \cmd{reverse\_proxy}. \\
Dépendance vulnérable & pip-audit, Trivy, SBOM &
Werkzeug 3.1.8 $\rightarrow$ 3.1.9 après alerte. \\
\bottomrule
\end{tabularx}
\end{center}

\begin{retenir}
Dix-huit menaces, dix-huit parades, et pour chacune un endroit où aller voir.
C'est ce dernier point qui compte : une mesure de sécurité qu'aucun test ne
verrouille est une mesure qui disparaîtra à la prochaine modification, sans que
personne s'en aperçoive.
\end{retenir}


\end{document}
